\section{Background}\label{sec:background}

Questo lavoro prende ispirazione da quello di
Csurilla\cite{csurillaLessObviousEffect2023}, il quale mostra come l'effetto
\textit{host} diventi meno significativo quando si tenga conto di variabili
aggiuntive, in particolare: PIL e popolazione (entrambi in logaritmo e aggregati
con media geometrica nei 4 anni precedenti), appartenenza passata al blocco
comunista (poiché tali Stati hanno storicamente propeso verso alti investimenti
nello sport \todo{citazione fonte del paper}), media di medaglie negli eventi
passati (stando ad indicare una predilezione per lo sport per motivi culturali),
\textit{dummy variable} separate per ogni anno (per assorbire avvenimenti
globali specifici e per tenere conto del numero variabile di medaglie
disponibili di anno in anno). Le variabili relative all'\textit{hosting},
invece, \textit{HOST<Year>}, \textit{PRE<Year>} e \textit{POST<Year}, sono
separate per ogni anno di Olimpiadi e valgono \textit{1}, rispettivamente,
quando un Paese ospita, quando ospiterà nella successiva edizione e quando ha
ospitato nella precedente; dunque, \textit{PRE} cerca di individuare un
miglioramento prestazionale generale cominciato alcuni anni prima di ospitare,
mentre \textit{POST} cerca di catturare un \textit{boost} duraturo fornito dal
giocare in casa.\\
A differenza di ricerche pregresse, il paper impiega una regressione di tipo
\textit{zero-inflated negative binomial} (ZINB), per ovviare ad alcune
limitazioni di altri modelli su questi dataset. Il numero di medaglie, infatti,
è una variabile discreta, con alta varianza e tante istanze pari a 0 (le Nazioni
partecipanti sono per la maggior parte molto piccole): altri modelli, come
\textit{regressione lineare}, \textit{Poisson regression} e
\textit{Tobit}\todo{refs}, soffrono per almeno uno di questi motivi. Le
peculiarità del dataset diventano ancora più acuite per il fatto che Csurilla
abbia applicato la regressione dividendo il dataset a livello di sport (oltre
che per anno e Paese).\\
\\
Nel corso del lavoro, denoteremo le variabili come indicato di seguito. Alcune
sono appena state descritte, tra cui \textit{HOST<Year>} (detto anche
\textit{OG<Year>}), \textit{PRE<Year>} e \textit{POST<Year>} (quando non è
specificato l'anno, si intende che è stata usata un'unica variabile per tutti
insieme). Si noti che mancherà \textit{PRE} per l'ultima edizione.
\textit{MACRO} è l'insieme delle variabili macroeconomiche aggiuntive:
\textit{GDP} per il PIL (GDPpc, o pro capite, in dollari del 2011, aggiustato per
l'inflazione), \textit{POP} per la popolazione (in migliaia), \textit{COMM} per
l'appartenenza al blocco comunista, \textit{AM} per la media storica di
medaaglie. \textit{YEAR<Year>} è una \textit{dummy variable} per ogni anno;
\textit{BOYCOTT<Year>}, in assenza di \textit{YEAR}, si ha solo per assorbire
sbilanciamenti nei due anni di boicottaggio reciproco tra USA e URSS (1980 e
1984): essendo questi i due Paesi più vincenti, la mancanza di uno ha fatto sì
che l'altro ottenesse la maggior parte delle medaglie, influenzando notevolmente
i dati. \textit{const} e \textit{alpha} sono variabili specifiche dei modelli.\\
\textit{CLOSE\_<group><year>} vale 1 per le Nazione appartenenti al dato gruppo
quando una di esse abbia ospitato i Giochi, ma 0 per chi ospita (senza
\textit{<year>} è aggregata per anno). I gruppi sono stati scelti per cercare di
catturare l'effetto di competere in un Paese vicino che non fosse il proprio,
andando così ad escludere i fattori geografici. L'Europa è apparso il luogo
ideale: massima differenza di 1 ora; cultura e ideali simili (soprattutto sotto
l'Unione Europea); economia, trasporti e alloggi di alto livello; gran numero
Paesi, soprattutto tra quelli più forti (dunque, alto numero di dati per
edizioni ospitate), e per di più piccoli (viaggiare in Europa può essere più
facile che da una parte all'altra della Cina). Dunque, abbiamo: Dunque, abbiamo:
\begin{itemize}
    \item \textit{CLOSE\_WEST} (Belgio, Francia, Paesi Bassi)
    \item \textit{CLOSE\_CENTER} (Austria, Repubblica Ceca, Germania, Germania
    Ovest, Ungheria, Polonia, Slovenia, Slovacchia, Svizzera)
    \item \textit{CLOSE\_GMT1} (\textit{CLOSE\_CENTER} + \textit{CLOSE\_WEST} +
    Italia + Spagna): fuso orario GMT+1
    \item \textit{CLOSE\_MAIN} (Francia, Germania, Italia, Spagna): le squadre
    più forti in Europa continentale
    \item \textit{CLOSE\_WIDE} (\textit{CLOSE\_GMT1} + Grecia + UK): gruppo
    allargato a due ospitanti importanti
\end{itemize}
Le suddivisioni Ovest e Centrale seguono i confini culturali e spaziali definiti
dalla StAGN\cite{CentralEurope2026}.\\
Inoltre, \textit{OLS} (Ordinary Least Squares) indica il metodo di regressione
lineare utilizzato in seguito.\\
Faremo riferimento ai team anche come Paesi; d'altronde, quelli senza uno Stato
non sono stati considerati, non avendo dati macroeconomici.\\
